%==============================================================================
%  PROJETO DE CIÊNCIA DE DADOS IV
%  Entrega 1 --- Problema de pesquisa e revisão bibliográfica
%
%  COMO USAR ESTE TEMPLATE
%  -----------------------
%  1. Preencha os dados do grupo no bloco "DADOS DO GRUPO", logo abaixo.
%  2. Escreva o conteúdo nas seções marcadas com  % >>> ESCREVA AQUI
%  3. Cadastre as referências no arquivo  referencias.bib
%  4. Compile (botão Recompile). O Overleaf já compila quantas vezes for
%     preciso para as citações aparecerem.
%
%  Linhas que começam com %  são comentários: não aparecem no PDF.
%==============================================================================
\documentclass[12pt,a4paper]{article}

%------------------------------------------------------------------------------
%  Pacotes
%------------------------------------------------------------------------------
\usepackage[T1]{fontenc}
\usepackage[utf8]{inputenc}
\usepackage[brazil]{babel}
\usepackage{lmodern}
\usepackage[a4paper,top=3cm,bottom=2.5cm,left=3cm,right=2.5cm]{geometry}
\usepackage{graphicx}
\usepackage{xcolor}
\usepackage{booktabs}
\usepackage{tabularx}
\usepackage{array}
\usepackage{indentfirst}
\usepackage{setspace}
\usepackage[hidelinks]{hyperref}

%  Citações e referências no padrão ABNT (autor-data).
%  O comando \cite{chave} produz, por exemplo: (SILVA; COSTA, 2024)
\usepackage[alf]{abntex2cite}

\definecolor{pucverde}{HTML}{00913F}

\onehalfspacing
\setlength{\parindent}{1.25cm}

\newcolumntype{Y}{>{\raggedright\arraybackslash}X}

%==============================================================================
%  DADOS DO GRUPO --- PREENCHA ESTE BLOCO
%==============================================================================
\newcommand{\TITULOPROJETO}{Desigualdade geográfica no acesso a modelos de IA generativa comercial}

\newcommand{\INTEGRANTEUM}{Pedro Dias Soares}
\newcommand{\INTEGRANTEDOIS}{Eduardo Velloso Landeiro}
\newcommand{\INTEGRANTETRES}{\textit{Composição do grupo a confirmar}}

\newcommand{\TEMADOGRUPO}{Desigualdade geográfica no acesso à IA generativa comercial}

%  Cole aqui o link de LEITURA do projeto no Overleaf.
%  (Share > Link sharing > "Anyone with this link can view" > copie o link)
%  Se o endereço contiver #, escreva \# entre as chaves abaixo.
\newcommand{\LINKPROJETO}{}
\newcommand{\MOSTRARLINKPROJETO}{%
  \if\relax\detokenize\expandafter{\LINKPROJETO}\relax
    \emph{Pendente de inserção e teste sem login}%
  \else
    \expandafter\url\expandafter{\LINKPROJETO}%
  \fi
}

%==============================================================================
\begin{document}
\hypersetup{pageanchor=false}

%------------------------------------------------------------------------------
%  CAPA
%------------------------------------------------------------------------------
\begin{titlepage}
\begin{center}

\IfFileExists{puclogo.pdf}%
  {\includegraphics[page=1,height=2.6cm]{puclogo.pdf}}%
  {}

\vspace{0.6cm}
{\large Pontifícia Universidade Católica de Minas Gerais\par}
{\normalsize Instituto de Ciências Exatas e Informática\par}
{\normalsize Curso de Ciência de Dados e Inteligência Artificial\par}

\vspace{2.5cm}
{\color{pucverde}\rule{0.6\linewidth}{1pt}}

\vspace{0.8cm}
{\LARGE\bfseries \TITULOPROJETO \par}

\vspace{0.6cm}
{\large Entrega 1 --- Problema de pesquisa e revisão bibliográfica\par}

\vspace{0.4cm}
{\color{pucverde}\rule{0.6\linewidth}{1pt}}

\vspace{2.5cm}
{\large
\INTEGRANTEUM\\[4pt]
\INTEGRANTEDOIS\\[4pt]
\INTEGRANTETRES\par}

\vspace{1.5cm}
{\normalsize\textbf{Tema:} \TEMADOGRUPO\par}

\vfill
{\small\textbf{Base editável para preenchimento pelo grupo}\par}
{\small\textbf{Projeto LaTeX:} \MOSTRARLINKPROJETO\par}

\vspace{0.6cm}
{\normalsize Belo Horizonte\\ 2026\par}

\end{center}
\end{titlepage}
\hypersetup{pageanchor=true}

%------------------------------------------------------------------------------
\section{Introdução}
%------------------------------------------------------------------------------
% >>> ESCREVA AQUI: de 2 a 4 parágrafos situando o tema.
%
% O que entra nesta seção:
%   - do que se trata o assunto;
%   - por que ele é relevante;
%   - qual o contexto em que o problema aparece.
%
% O que NÃO entra: a pergunta de pesquisa (ela vem na Seção 2) nem a
% discussão dos trabalhos encontrados (ela vem na Seção 3).
%
% Para citar um trabalho, use \cite{chave}, onde "chave" é o identificador
% que está no arquivo referencias.bib. Exemplos:
%   Segundo \citeonline{silva2024exemplo}, o problema ...   ->  Silva e Costa (2024)
%   O problema é recorrente na área \cite{silva2024exemplo}. ->  (SILVA; COSTA, 2024)

Escreva aqui a introdução.

%------------------------------------------------------------------------------
\section{Problema de pesquisa}
%------------------------------------------------------------------------------

\subsection{Problema}
% >>> ESCREVA AQUI: um parágrafo enunciando a dificuldade ou lacuna.
%
% O problema aponta uma tensão, mas ainda não é uma pergunta. Exemplo:
% "Modelos treinados em bases estrangeiras podem não funcionar em dados
%  brasileiros, e não se sabe o tamanho dessa perda de desempenho."

Escreva aqui o problema.

\subsection{Pergunta de pesquisa}
% >>> ESCREVA AQUI: UMA frase, terminando em ponto de interrogação.
%
% Teste antes de escrever: qual seria uma resposta possível? Se você consegue
% imaginar duas ou três respostas diferentes e plausíveis, é pergunta.
% Se a única resposta que ocorre é "depende", ainda é tema.

Escreva aqui a pergunta, terminando com ponto de interrogação?

\subsection{Objetivo geral}
% >>> ESCREVA AQUI: uma frase, começando com verbo no infinitivo,
% correspondendo diretamente à pergunta acima.

Escreva aqui o objetivo geral.

\subsection{Objetivos específicos}
% >>> ESCREVA AQUI: de 3 a 5 itens.
%
% Para cada um, pergunte: "como eu saberei que terminei este?"
% Se não houver resposta objetiva, troque o verbo.
%
% Verbos que funcionam:     levantar, caracterizar, implementar, comparar,
%                           avaliar, medir, identificar, classificar
% Verbos problemáticos:     compreender, estudar, explorar, conhecer, abordar

\begin{enumerate}
  \item Primeiro objetivo específico.
  \item Segundo objetivo específico.
  \item Terceiro objetivo específico.
\end{enumerate}

\subsection{Justificativa}
% >>> ESCREVA AQUI: até 5 linhas.
%
% Apoie-se em pelo menos uma destas fontes:
%   - a lacuna na literatura (o que a Seção 3 vai mostrar);
%   - a consequência prática (alguém decidiria diferente sabendo a resposta);
%   - o risco de errar (aplicar o método errado tem custo mensurável).
%
% Teste: se o texto serviria para qualquer outro trabalho da área, reescreva.

Escreva aqui a justificativa.

\subsection{Delimitação de escopo}
% >>> ESCREVA AQUI: no mínimo 3 itens do que o trabalho NÃO vai fazer.

Este trabalho não vai tratar de:

\begin{itemize}
  \item Primeiro item fora do escopo.
  \item Segundo item fora do escopo.
  \item Terceiro item fora do escopo.
\end{itemize}

%------------------------------------------------------------------------------
\section{Revisão bibliográfica}
%------------------------------------------------------------------------------
% >>> ESCREVA AQUI: de 2 a 3 páginas.
%
% A REGRA DO PARÁGRAFO
% Cada parágrafo defende UMA afirmação; os trabalhos citados são a evidência
% dela. A afirmação vai na primeira frase; os autores entram depois.
%
% Assim NÃO:
%   "Silva e Costa (2024) propuseram uma arquitetura baseada em LSTM e
%    obtiveram F1 de 0,87. Já Oliveira (2022) aplicou gradient boosting..."
%   (isso é uma lista com conectivos)
%
% Assim SIM:
%   "A escolha da métrica domina a comparação entre métodos nesta tarefa.
%    Três dos trabalhos analisados relatam vantagem de modelos profundos sob
%    acurácia, mas dois deles invertem a conclusão quando reportam PR-AUC
%    \cite{silva2024exemplo,oliveira2022exemplo}. Isso sugere que diferenças
%    publicadas podem refletir a escolha de métrica tanto quanto a capacidade
%    do modelo."
%
% ESTRUTURA
% Organize por DIMENSÃO do problema (os aspectos em que os trabalhos
% divergem), não por ordem cronológica nem um parágrafo por artigo.
% As dimensões saem da matriz de síntese: são as colunas em que os
% trabalhos mais discordam entre si.
%
% O ÚLTIMO PARÁGRAFO
% Enuncia a lacuna. Ela deve ser CONSEQUÊNCIA do que veio antes.
% Teste: cubra o último parágrafo. A lacuna ainda é previsível? Se não,
% ela foi colada ali.

Escreva aqui a revisão bibliográfica.

%------------------------------------------------------------------------------
%  REFERÊNCIAS --- geradas automaticamente do arquivo referencias.bib
%------------------------------------------------------------------------------
\bibliography{referencias}

%------------------------------------------------------------------------------
\appendix
%------------------------------------------------------------------------------

\section{Protocolo de busca}
% >>> PREENCHA a tabela com as buscas que vocês fizeram.
% Mínimo de duas bases, sendo pelo menos uma biblioteca digital científica
% (ACM DL, IEEE Xplore, SBC OpenLib) ou base indexadora (Scopus, Web of Science).

\vspace{6pt}
{\footnotesize
\noindent\begin{tabularx}{\linewidth}{@{}p{2.1cm}p{1.8cm}Yp{2.0cm}p{1.1cm}@{}}
\toprule
\textbf{Base} & \textbf{Data} & \textbf{Expressão de busca} & \textbf{Campos} & \textbf{Total} \\
\midrule
 & & & & \\[10pt]
\midrule
 & & & & \\[10pt]
\midrule
 & & & & \\[10pt]
\bottomrule
\end{tabularx}}

\vspace{12pt}
\subsection*{Critérios de inclusão}
% >>> LISTE os critérios. Bom critério é verificável olhando o artigo.
\begin{itemize}
  \item Primeiro critério de inclusão.
  \item Segundo critério de inclusão.
\end{itemize}

\subsection*{Critérios de exclusão}
\begin{itemize}
  \item Primeiro critério de exclusão.
  \item Segundo critério de exclusão.
\end{itemize}

\newpage
\section{Matriz de síntese}
% >>> COPIE o bloco abaixo uma vez para CADA trabalho incluído (6 a 8 no total).
%
% Campo que o artigo não permite preencher recebe a marca [não informado].
% Essa marca é um achado, não uma falha de leitura.
%
% Coloque as INICIAIS de quem preencheu cada trabalho.

\vspace{6pt}
{\footnotesize
\noindent\textbf{Trabalho 1} \hfill \textit{Preenchido por:} \underline{\hspace{2.5cm}}

\vspace{2pt}
{\renewcommand{\arraystretch}{2.0}
\noindent\begin{tabularx}{\linewidth}{@{}p{2.4cm}Y@{}}
\toprule
\textbf{Referência} & \\
\textbf{Problema} & \\
\textbf{Dados} & \\
\textbf{Método} & \\
\textbf{Baseline} & \\
\textbf{Métrica} & \\
\textbf{Resultado} & \\
\textbf{Limitação} & \\
\bottomrule
\end{tabularx}}

\vspace{14pt}
\noindent\textbf{Trabalho 2} \hfill \textit{Preenchido por:} \underline{\hspace{2.5cm}}

\vspace{2pt}
{\renewcommand{\arraystretch}{2.0}
\noindent\begin{tabularx}{\linewidth}{@{}p{2.4cm}Y@{}}
\toprule
\textbf{Referência} & \\
\textbf{Problema} & \\
\textbf{Dados} & \\
\textbf{Método} & \\
\textbf{Baseline} & \\
\textbf{Métrica} & \\
\textbf{Resultado} & \\
\textbf{Limitação} & \\
\bottomrule
\end{tabularx}}
}

\newpage
\section{Nota de autoria}

% >>> PREENCHA com os nomes reais.
% Lembre: quem confere não escreve, e quem escreve não confere o próprio texto.

\begin{description}
  \item[Matriz de síntese preenchida por:] \hfill \\
        Nome (trabalhos 1 a 3), Nome (4 a 6), Nome (7 e 8)

  \item[Texto da revisão bibliográfica redigido por:] \hfill \\
        Nome

  \item[Conferência de fidelidade] --- cada afirmação do texto foi checada
        contra a fonte citada --- feita por: \hfill \\
        Nome

  \item[Conferência de argumento] --- verificação de que cada parágrafo defende
        uma ideia e de que a lacuna decorre do texto --- feita por: \hfill \\
        Nome
\end{description}

\end{document}
